% ===================================================================
\section{Inicio}

\captura{inicio_1.png}

La página de inicio resume el modelo en una tabla. Cada fila es una capa del modelo,
lo que modela y la página donde se explora.

\begin{description}
\item[\boton{Cargar escenario de ejemplo “Bre-B estilizado”}.] Reemplaza la red y los
parámetros por un escenario de ejemplo: red núcleo-periferia de 245 entidades (del orden
de los participantes de Bre-B en su primer año), grado medio 8, índice de amplificación
$R=0.85$, redistribución $D_0=0.05$, acople $K=3$ y colchón $0.5$. Los valores son
ilustrativos.
\item[\boton{Volver a la red por defecto}.] Restaura la red inicial: núcleo-periferia de
80 entidades, grado medio 6, $R=0.8$, semilla 42.
\end{description}

Debajo, cuatro métricas describen la red activa:

\begin{description}
\item[Entidades.] Número de nodos.
\item[Aristas.] Número de relaciones de pago dirigidas.
\item[Grado medio.] Número medio de pagadores (o de receptores) por entidad.
\item[CV de montos.] Coeficiente de variación de los montos: desviación estándar
dividida por la media. Un valor alto indica pocos pagos muy grandes y muchos pequeños.
\end{description}

% ===================================================================
\section{Construir red}

\captura{construir_1.png}

\begin{representa}
Un sistema de pagos como red. Cada círculo es una entidad y cada línea un pago
recurrente. El tamaño del círculo es proporcional al volumen de pagos. En la topología
núcleo-periferia, los círculos rojos son el núcleo: las grandes entidades que concentran
los pagos.
\end{representa}

\subsection*{Controles (dentro del formulario)}
\begin{controles}
Topología & Erdős–Rényi, Barabási–Albert, Núcleo-periferia (núcleo-periferia) &
Forma de la red. \textbf{Erdős–Rényi}: conexiones al azar, todas las entidades parecidas.
\textbf{Barabási–Albert}: unas pocas entidades muy conectadas (concentradores).
\textbf{Núcleo-periferia}: un núcleo denso y una periferia conectada a él; es la forma de
los sistemas de pagos reales. \\
Número de entidades & 10 a 600 (80) & Tamaño de la red. Por encima de 300 se desactiva la
simulación en vivo. \\
Grado medio & 1 a 20 (6) & Número medio de contrapartes de cada entidad. Más grado, red
más densa. Debe ser menor que el número de entidades menos uno. \\
Dispersión de montos ($\sigma$ lognormal) & 0 a 2 (0.8) & Desigualdad de los montos. Con
0 todos los pagos valen igual; con valores altos aparecen pagos muy grandes. \\
Fracción en el núcleo & 0.02 a 0.40 (0.10) & Proporción de entidades del núcleo. Solo se
activa con núcleo-periferia. \\
Semilla & 0 a $10^6$ (42) & Fija el azar de la construcción. Cambiarla produce otra red
con las mismas características. \\
\end{controles}

Al presionar \boton{Construir} se genera la red. Las simulaciones previas se borran
porque correspondían a la red anterior.

\subsection*{Resultados}
\begin{description}
\item[Entidades y aristas dirigidas.] Tamaño de la red.
\item[Densidad.] Fracción de las conexiones posibles que existen, entre 0 y 1.
\item[CV de montos.] Ver \emph{Inicio}.
\item[Gráfica de la red.] Al pasar el cursor sobre una entidad se ve cuánto paga y
cuánto recibe por ciclo.
\end{description}

\begin{pruebe}
Construya una red Barabási–Albert de 100 entidades y luego una Erdős–Rényi del mismo
tamaño y grado. En la primera verá unos pocos círculos grandes en el centro; en la
segunda, círculos de tamaño parecido. Después vaya a \emph{Diagnóstico}: con el mismo
$\beta$, la red Barabási–Albert tendrá un $R$ mayor.
\end{pruebe}

% ===================================================================
\section{Diagnóstico}

\captura{diagnostico_1.png}

\begin{representa}
La salud de la red frente a un choque, calculada sin simular, a partir de los valores
propios de la red. Responde: si una entidad sufre un déficit, ¿la red lo absorbe o lo
amplifica? ¿Cuánto tarda en recuperarse? ¿Qué entidades son más importantes?
\end{representa}

\subsection*{Controles}
\begin{controles}
$\alpha$ · autocorrección [1/ciclo] & 0.05 a 3 (1.0) & Rapidez con que cada entidad
corrige su déficit por sí misma. $1/\alpha$ es el tiempo típico de corrección. Subirlo
estabiliza la red. \\
$\beta$ · amplificación & 0 a $2\alpha/\rho$ (el que da $R=0.8$) & Déficit que genera en
una entidad el déficit de quien le paga. El rango se ajusta a la red para que $R$ vaya de
0 a 2. Subirlo desestabiliza la red. \\
$D_0$ · redistribución & 0 a 2 (0) & Traslado de presión de liquidez de una entidad a sus
receptores sin crearla ni destruirla. Acelera la recuperación. \\
Comunidades & 2 a 8 (3) & Número de grupos en que se divide la red para colorearla. \\
\end{controles}

\subsection*{Resultados}
\begin{description}
\item[Índice de amplificación $R$.] $R=\beta\rho(W)/\alpha$, donde $\rho(W)$ es el radio
espectral de la red (su valor propio dominante). $R<1$: la red amortigua. $R>1$: la red
amplifica. Es el análogo del número reproductivo de una epidemia.
\item[Estabilidad $\sigma$ con semáforo.] Mayor parte real de los valores propios de
$M=-\alpha I-D_0L+\beta W$. Incluye la redistribución, así que es más completo que $R$.
\textcolor{verde}{Verde}: $\sigma<-0.05$, la red amortigua.
\textcolor{ambar}{Amarillo}: $|\sigma|\le0.05$, la red está cerca del límite y los
choques tardan mucho en disiparse. \textcolor{rojo}{Rojo}: $\sigma>0.05$, el déficit
crece sin límite.
\item[Tiempo de recuperación.] $\tau_\text{rec}=-1/\sigma$, en ciclos: tiempo en que un
choque se reduce a un 37~\% ($1/e$). Aparece $\infty$ si la red no se recupera.
\item[Conectividad algebraica $\lambda_2$.] Mide qué tan bien conectada está la red como
un todo. Vale 0 si está partida en pedazos aislados. Mayor $\lambda_2$, más rápida la
redistribución.
\item[Comunidades.] Gráfica con la red coloreada por grupos que se pagan mucho entre sí y
poco con el resto. Se obtienen con el vector de Fiedler (bisección espectral).
\item[Ranking por vector propio dominante.] Las 15 entidades más importantes.
\emph{importancia}: cuánto déficit propaga si entra en problemas (vector propio
izquierdo). \emph{vulnerabilidad}: cuánto déficit acumula si otros entran en problemas
(vector propio derecho). \emph{paga} y \emph{recibe}: flujos por ciclo.
\end{description}

\captura{diagnostico_2.png}

\subsection*{Choque de prueba (sin impagos)}
Aplica un retiro de liquidez a una entidad y muestra cómo evoluciona el déficit, sin
umbrales: nadie entra en impago. Sirve para ver la estabilidad en acción.

\begin{controles}
Entidad chocada & 0 a $n-1$ (la más importante) & Entidad que sufre el choque. \\
Tamaño & 0.1 a 10 ciclos (2) & Déficit inicial, medido en ciclos de los pagos salientes
de la entidad: 2 equivale a dos ciclos de lo que paga. \\
\end{controles}

\begin{description}
\item[Trayectorias.] El déficit de cada entidad en el tiempo; la entidad chocada aparece
resaltada.
\item[Déficit total del sistema ($\Sigma x$).] Suma de todos los déficits. Con $R<1$
termina en cero; con $R>1$ crece.
\end{description}

\begin{pruebe}
Mueva $\beta$ hasta que $R$ pase de 0.9 a 1.1. El semáforo cambia de verde a rojo y el
déficit total del choque de prueba deja de bajar y empieza a crecer. Luego suba $D_0$:
$\sigma$ baja y la red puede volver a ser estable aunque $R>1$.
\end{pruebe}

% ===================================================================
\section{Ritmo de pagos}

\captura{ritmo_1.png}

\begin{representa}
Cada entidad paga en ciclos, como un reloj que da vueltas. La posición en la vuelta es
su fase. Cada entidad tiene su propio ritmo, pero tiende a ajustarlo al de quienes le
pagan. La página muestra si las entidades terminan sincronizadas y cuánto déficit se
genera cuando el ingreso llega desfasado.
\end{representa}

\subsection*{Controles}
\begin{controles}
$K$ · dependencia de fondos entrantes & 0 a 10 (2) & Intensidad con que cada entidad
ajusta su ritmo al de quienes le pagan. Más $K$, más sincronía. \\
$\Delta$ · dispersión de ritmos naturales & 0 a 2 (0.5) & Qué tan distintos son los
ritmos propios de las entidades. Más $\Delta$, más difícil sincronizar. \\
$\gamma$ · peso de la desincronización en el déficit & 0 a 0.5 (0) & Cuánto déficit genera
el ingreso desfasado. Con 0, la sincronía no afecta la liquidez. \\
Desconectar entidad a mitad de la corrida & ninguna o una de las 30 que más pagan &
Elimina todos los pagos de esa entidad en la mitad de la simulación (ciclo 10 de 20). \\
\end{controles}

\subsection*{Resultados}
\begin{description}
\item[$r$ final.] Parámetro de orden al final: 0 si las fases están dispersas, 1 si todas
coinciden. Con una desconexión, la flecha muestra la diferencia frente a la red
completa.
\item[Déficit total final.] Suma de los déficits al final. Con desconexión, muestra la
diferencia.
\item[Coherencia local media.] Promedio de $r_i$, que mide si a cada entidad le llega el
ingreso en fase (1), sin relación (0) o en contrafase ($-1$).
\item[Círculo de fases.] Animación: cada punto es una entidad sobre un círculo según su
fase. Si se agrupan, están sincronizadas. Use el botón de reproducción.
\item[Parámetro de orden $r(t)$.] Evolución de la sincronía. Con desconexión se muestran
dos curvas: con y sin la entidad.
\item[Déficit total.] Evolución de la suma de déficits.
\item[Coherencia local al final.] Red coloreada por $r_i$: verde, ingreso en fase; rojo,
en contrafase.
\end{description}

\captura[0.8]{ritmo_2.png}

\begin{aviso}
Con $\gamma=0$ aparece un mensaje: las fases se sincronizan, pero el déficit no cambia.
Suba $\gamma$ para ver el efecto del ritmo sobre la liquidez.
\end{aviso}

\begin{pruebe}
Con $\gamma=0.1$, mueva $K$ de 0 a 6. Con $K$ bajo, $r$ queda cerca de 0 y el déficit
total crece; con $K$ alto, $r$ se acerca a 1 y el déficit se reduce. Después desconecte la
primera entidad de la lista: $r$ cae y el déficit sube en la segunda mitad.
\end{pruebe}

% ===================================================================
\section{Simulador de choques}

\captura{simulador_1.png}

\begin{representa}
Un choque sobre una entidad y la cascada de impagos que provoca. A diferencia del
choque de prueba de \emph{Diagnóstico}, aquí sí hay colchones: la entidad que los supera
deja de pagar y sus receptores pierden ingreso. Usa los parámetros $\alpha$, $\beta$,
$D_0$, $K$ y $\gamma$ elegidos en las páginas anteriores.
\end{representa}

\subsection*{Controles (dentro del formulario)}
\begin{controles}
Entidad chocada & todas, ordenadas por lo que pagan & Entidad que recibe el choque. Junto
al número aparece cuánto paga por ciclo. \\
Tipo de choque & Retiro de liquidez o Caída operativa & \textbf{Retiro de liquidez}: un
salto súbito en el déficit, como una retirada masiva de depósitos. \textbf{Caída
operativa}: la entidad sigue recibiendo pero no paga durante un tiempo, como una falla
tecnológica. \\
Tamaño del retiro & 0.1 a 10 ciclos (2) & Solo para el retiro: déficit inicial en ciclos
de pagos salientes. \\
Duración de la caída $T$ & 0.5 a 15 ciclos (3) & Solo para la caída operativa: cuánto
dura sin pagar. \\
Colchón $c$ & 0.05 a 3 (0.5) & Reserva de cada entidad, medida en ciclos de sus pagos
salientes. Más colchón, menos impagos. \\
$\lambda$ · pérdida por ingreso no recibido & 0 a 3 (1) & Cuánto déficit genera perder un
pagador. Con 1, perder todo el ingreso equivale a un ciclo de los pagos propios. \\
\end{controles}

El choque ocurre en el ciclo 1 y la simulación dura 20 ciclos.

\subsection*{Resultados}
\begin{description}
\item[En impago.] Porcentaje de entidades que terminan en impago. El símbolo \textbf{?}
muestra el número.
\item[Liquidez no pagada.] Pagos que no se hicieron por los impagos, acumulados, en pagos
medios.
\item[DebtRank.] Medida clásica del daño sistémico (Battiston et al., 2012), entre 0 y
1: fracción del valor económico de la red afectada por el estrés, sin contar a la
entidad chocada.
\item[Primer contagio.] Ciclo en que cae la primera entidad distinta de la chocada.
``—'' si no hubo contagio.
\item[Colapso global.] Mensaje rojo si cae el 30~\% o más de la red.
\item[Animación de la cascada.] Red coloreada por estado: \textcolor{verde}{verde}, sin
problemas; \textcolor{ambar}{ámbar}, en estrés (déficit mayor que la mitad de su
colchón); \textcolor{rojo}{rojo}, en impago. La entidad chocada tiene borde negro. Use el
botón de reproducción o arrastre la barra de tiempo.
\item[Línea de tiempo de eventos.] Cada punto es un impago: en el eje horizontal el
ciclo y en el vertical la entidad.
\item[Liquidez no pagada acumulada.] Evolución en el tiempo del indicador anterior.
\item[\boton{Exportar indicadores y eventos (CSV)}.] Descarga los resultados en un
archivo que se abre con Excel.
\end{description}

\begin{aviso}
Si $\gamma>0$ algunas entidades pueden caer antes del choque por la desincronización.
Un mensaje lo indica; esas caídas no cuentan como contagio.
\end{aviso}

\captura{simulador_2.png}

\subsection*{Entidades más peligrosas}
El botón \boton{Calcular ranking} choca, una por una, cada entidad con el tamaño elegido
y mide la fracción media de las demás que caen (\emph{cascada\_media}). Muestra las 10 más
peligrosas. Tarda unos segundos y el resultado queda guardado mientras no cambie la
configuración.

\begin{pruebe}
Choque la entidad que más paga con un retiro de 3 ciclos y colchón 0.5. Anote el
porcentaje en impago. Suba el colchón a 1.0 y repita: la cascada se reduce mucho.
Luego compare un retiro con una caída operativa de 3 ciclos sobre la misma entidad.
\end{pruebe}

% ===================================================================
\section{Explorar escenarios}

\captura{explorar_1.png}

\begin{representa}
Los resultados de 20\,000 simulaciones hechas de antemano sobre 300 redes distintas,
con parámetros variados. Permite responder preguntas generales sin simular: ¿en qué
condiciones colapsa una red? ¿Qué topología es más frágil?
\end{representa}

\begin{aviso}
Requiere el conjunto de datos. Si no existe, la página indica el comando para generarlo:
\texttt{uv run pulso sweep configs/sweep.yaml} (unos 20 minutos).
\end{aviso}

\subsection*{Filtros (en el desplegable \emph{Filtros})}
\begin{controles}
Topología & las tres (todas) & Redes que se incluyen. \\
Tipo de choque & liquidez, operativo (ambos) & Choques que se incluyen. \\
Índice de amplificación $R$ & 0.4 a 1.6 (todo) & Rango de $R$. \\
Grado medio & 2 a 12 (todo) & Rango de conectividad. \\
Colchón $c$ & 0.2 a 1.5 (todo) & Rango de colchones. \\
Acople $K$ & 0 a 6 (todo) & Rango de acople de fase. \\
\end{controles}

\subsection*{Resultados}
\begin{description}
\item[Corridas filtradas.] Número de simulaciones que cumplen los filtros.
\item[P(colapso).] Porcentaje de ellas con colapso global (30~\% o más en impago).
\item[Fracción media en impago.] Promedio del porcentaje de entidades que caen.
\item[Liquidez no pagada media.] En ciclos de flujo total del sistema.
\item[Ventana de contagio.] Fracción en impago contra grado medio, por topología. Muestra
que las cascadas son mayores con conectividad intermedia: con pocas conexiones el choque
no se propaga; con muchas, cada entidad diversifica su ingreso y absorbe la pérdida.
\item[Fracción en impago contra $R$.] Cada punto es una simulación (máximo 3000). Por
encima de $R=1$ aparecen las cascadas grandes.
\item[Distribución del tamaño de cascada.] Histograma por tipo de choque. La mayoría de
los choques se absorben; unos pocos producen colapsos.
\end{description}

\captura{explorar_2.png}

\subsection*{Comparar escenarios lado a lado}
Elija hasta cuatro escenarios por su número (por defecto, los dos con mayor cascada). La
tabla muestra sus parámetros y resultados en columnas. \boton{Exportar filtro (CSV)}
descarga todos los escenarios filtrados.

\begin{pruebe}
Filtre $R$ entre 0.4 y 0.9 y observe P(colapso). Luego filtre $R$ entre 1.1 y 1.6. La
probabilidad de colapso aumenta mucho. Compare también las tres topologías por separado.
\end{pruebe}

% ===================================================================
\section{Mapa de riesgo}

\captura{mapa_1.png}

\begin{representa}
El riesgo de cada entidad de la red activa, predicho por modelos de aprendizaje
automático entrenados con las 20\,000 simulaciones. La predicción tarda milisegundos;
simular cada entidad tardaría minutos. El riesgo de una entidad es el daño que causaría si
sufriera un retiro de 3 ciclos.
\end{representa}

\begin{aviso}
Requiere los modelos entrenados. Si no existen, la página indica el comando:
\texttt{uv run pulso train configs/train.yaml}.
\end{aviso}

\subsection*{Controles}
\begin{controles}
Modelo & Cascada (GCN), Cascada (XGBoost), P(colapso global) & Qué riesgo colorea la red.
\textbf{Cascada}: fracción de las demás entidades que caerían. La \textbf{GCN} es una red
neuronal que lee la estructura del grafo; \textbf{XGBoost} es un conjunto de árboles de
decisión. \textbf{P(colapso global)}: probabilidad de que caiga el 30~\% o más de la red.
\\
Entidad (sección SHAP) & todas, de mayor a menor riesgo & Entidad cuya predicción se
explica. \\
\end{controles}

\subsection*{Resultados}
\begin{description}
\item[Riesgo predicho por entidad.] Red coloreada de amarillo (bajo) a rojo (alto).
\item[Tabla.] Las 12 entidades más riesgosas con las tres predicciones y lo que pagan.
\item[Contribuciones SHAP.] Barras que explican la predicción de XGBoost para la entidad
elegida. Parten de un valor base (el riesgo promedio) y cada barra suma o resta riesgo:
\textcolor{rojo}{rojo} aumenta, \textcolor{azul}{azul} reduce. Los nombres de las variables
se explican en la tabla siguiente.
\item[\boton{Verificar con simulación}.] Simula de verdad el choque sobre la entidad y
compara el resultado (\emph{Simulado}) con las predicciones de la GCN y de XGBoost. La
flecha muestra el error de cada modelo.
\item[Métricas de validación.] Desplegable con el desempeño de los modelos fuera de
muestra.
\end{description}

\subsection*{Variables del gráfico SHAP}
\begin{center}
\small
\begin{longtable}{p{0.24\textwidth}p{0.68\textwidth}}
\toprule
\textbf{Variable} & \textbf{Significado} \\
\midrule
\endhead
\texttt{in\_degree}, \texttt{out\_degree} & Número de pagadores y de receptores. \\
\texttt{in\_amount}, \texttt{out\_amount} & Ingreso y pagos por ciclo. \\
\texttt{importance}, \texttt{vulnerability} & Importancia y vulnerabilidad espectrales
(ver \emph{Diagnóstico}). \\
\texttt{pagerank} & Centralidad PageRank. \\
\texttt{r\_local} & Coherencia local de fase antes del choque. \\
\texttt{threshold} & Colchón de la entidad. \\
\texttt{out\_share} & Fracción de todos los pagos del sistema que hace la entidad. \\
\texttt{shock\_to\_buffer} & Tamaño del choque dividido por el colchón propio. Mayor que 1:
la entidad cae con el choque. \\
\texttt{max\_dependency} & Mayor fracción del ingreso que algún receptor obtiene de esta
entidad. \\
\texttt{downstream\_impact} & Suma, sobre los receptores, de la pérdida que sufrirían
dividida por su colchón. \\
\texttt{first\_order\_defaults} & Fracción de entidades que caerían solo por la primera
ronda de pérdidas. \\
\texttt{debtrank} & DebtRank de un choque en la entidad. \\
\texttt{n}, \texttt{topology}, \texttt{R}, \texttt{sigma}, \texttt{K}, \texttt{gamma},
\texttt{buffer}, \texttt{loss} & Características de toda la red y del modelo: tamaño,
topología, amplificación, estabilidad, acople, desfase, colchón y pérdida. \\
\bottomrule
\end{longtable}
\end{center}

\begin{aviso}
Las predicciones solo son confiables dentro de los rangos con que se entrenaron los
modelos (los del archivo \texttt{configs/sweep.yaml}). El orden de las entidades dentro de
una misma red es aproximado; para una decisión, use \boton{Verificar con simulación}.
\end{aviso}

\begin{pruebe}
Elija la entidad más riesgosa, lea sus barras SHAP y presione \boton{Verificar con
simulación}. Compare el valor simulado con las dos predicciones.
\end{pruebe}

% ===================================================================
\section{Acerca del modelo}

\captura{acerca_1.png}

Página de consulta, sin controles. Contiene:

\begin{itemize}
\item Las ecuaciones del modelo de las páginas 1 a 7: déficit, ritmo de pagos e impago.
\item Una tabla con el significado de cada símbolo.
\item La unidad de tiempo, el integrador y la definición de los dos tipos de choque.
\item Los supuestos y límites: parámetros ilustrativos, redes sintéticas, acople en un
solo sentido, modelo lineal hasta el umbral y validez de las predicciones solo dentro de
los rangos de entrenamiento.
\item Las ecuaciones de la tarea 6 (páginas 8 y 9), el control nulo y las referencias.
\end{itemize}

% ===================================================================
\section{Cartera multicapa}

\captura{cartera_1.png}

\begin{representa}
Los clientes de un banco: personas (círculos) y empresas (cuadrados). Están unidos por
tres tipos de pago: \textbf{nómina} (empresa a empleado), \textbf{proveedores} (empresa a
empresa) y \textbf{transferencias} (persona a persona). Cada cliente tiene un ciclo de
pagos. Si su ingreso le llega desfasado de su propio ciclo, acumula déficit.

El banco no ve los ciclos: solo ve las transacciones. La página muestra el modelo
interno y lo que se estima a partir de las transacciones. El tiempo se mide en meses.
\end{representa}

\subsection*{Controles (dentro del formulario)}
\begin{controles}
Personas & 100 a 600 (400) & Número de personas de la cartera. \\
Firmas & 5 a 40 (20) & Número de empresas. \\
Informalidad & 0 a 0.9 (0.55) & Fracción de personas sin empleador en el banco. Los
informales tienen ritmos de pago más irregulares. \\
Semilla & 0 a 10\,000 (7) & Fija el azar: misma semilla, misma cartera. \\
Mundo & con red / nulo & \textbf{Con red}: la red transmite déficit (contagio, desfase y
pérdida de ingreso). \textbf{Nulo}: control; los pagos existen pero la red no crea
déficit. Mismos clientes y mismo azar. \\
$K$ nómina & 0 a 3 (0.8) & Cuánto ajusta un empleado su ritmo de pagos al de su
empleador. \\
$\gamma$ & 0 a 2 (0.5) & Cuánto déficit genera el ingreso desfasado. Con 0, la sincronía
no afecta el riesgo. \\
\end{controles}

\subsection*{Resultados principales}
\begin{description}
\item[Clientes.] Personas más empresas.
\item[$R=\rho(B)/\alpha$.] Índice de amplificación de la cartera. Menor que 1: la red
amortigua. En el mundo nulo vale 0.
\item[Impago en observación.] Porcentaje de clientes que caen en los primeros 6 meses
(ventana de observación).
\item[Impago a 12 meses.] Entre los que siguen activos a los 6 meses, porcentaje que cae en
los 12 meses siguientes. Es la etiqueta que predicen los modelos.
\item[Sincronía final $\hat r$.] Sincronía global de toda la cartera estimada desde las
transacciones. Suele ser baja porque las empresas pagan en fechas distintas: lo que
importa para el riesgo es la coherencia de cada cliente con sus propios pagadores, no la
de toda la cartera.
\item[Gráfica de la cartera.] El selector \emph{Color} cambia lo que se pinta:
\emph{coherencia local $\hat r$} (verde, ingreso en fase; rojo, desfasado),
\emph{impago a 12 meses} (rojo, cayó) o \emph{tipo de cliente}. Las líneas tienen el color
de su capa.
\item[Capas.] Tabla con el número de relaciones y el monto mensual de cada capa, en
salarios mínimos (SMMLV) por mes.
\item[Sincronía global de la cartera.] Curva gris: sincronía real del modelo, que el
banco no ve. Curva roja: la estimada desde las transacciones. La línea punteada marca el
fin de la observación (mes 6); la curva roja solo existe antes.
\end{description}

\captura{cartera_2.png}

\subsection*{¿Quién está desincronizado?}
\begin{description}
\item[Coherencia local $\hat r$ por grupo.] Diagramas de caja para empresas, personas
formales e informales, solo con clientes que reciben ingreso por la red. La caja abarca la
mitad central de los clientes y la línea es la mediana. Los formales están cerca de 1
(sincronizados con su nómina); los informales, cerca de 0.
\item[Impago a 12 meses según $\hat r$.] Tasa de impago por tramos de coherencia, con los
resultados del experimento completo de 20 carteras. En rojo el mundo con red; en gris el
nulo. El número sobre cada barra es la cantidad de clientes.
\end{description}

\begin{aviso}
En el mundo nulo también hay más impago con $\hat r$ bajo, porque los informales son más
riesgosos y están menos sincronizados. Por eso esta gráfica no prueba la hipótesis. La
prueba está en la comparación de modelos de la página siguiente.
\end{aviso}

\subsection*{Escenario: quiebra de un empleador}
\begin{controles}
Firma que quiebra en $t=1$ mes & las 8 con más empleados & Empresa que deja de pagar todo
(nómina, proveedores y cuota) en el primer mes. Se muestra su número de empleados en el
banco. \\
\end{controles}
Al presionar \boton{Simular quiebra} aparecen:
\begin{description}
\item[Empleados que caen (6 meses).] Porcentaje de sus empleados en impago entre el mes 1
y el 7, con la diferencia frente a la misma cartera sin quiebra.
\item[Resto de personas.] Mismo porcentaje para las demás personas.
\item[Empleados afectados.] Número de empleados de la firma.
\item[Clientes en impago acumulados.] Curvas con y sin quiebra.
\end{description}
En el mundo nulo la quiebra no transmite déficit y un mensaje lo recuerda.

\subsection*{Registro de transacciones}
El desplegable muestra las primeras 500 transacciones: lo que el banco ve. Columnas:
\texttt{t} (mes), \texttt{src} (pagador), \texttt{dst} (receptor; $-1$ es el banco, es
decir, una cuota), \texttt{layer} y \texttt{capa} (tipo de pago), \texttt{amount} (monto en
SMMLV). El botón descarga todas en CSV.

\begin{pruebe}
Con la cartera por defecto, simule la quiebra de la firma más grande en el mundo con red
y anote el porcentaje de empleados que caen. Cambie a mundo nulo y repita: la diferencia
desaparece. Después ponga $\gamma=0$ en el mundo con red y observe el color por $\hat r$:
la sincronía sigue igual, pero deja de afectar el impago.
\end{pruebe}

% ===================================================================
\section{Riesgo crediticio híbrido}

\captura{hibrido_1.png}

\begin{representa}
La respuesta a la pregunta del proyecto: ¿la sincronía de los ciclos de pago predice
impagos que el score clásico y la red estática no ven? Compara tres niveles de
variables con el mismo modelo y las mismas particiones:
\begin{enumerate}
\item \textbf{Clásico}: ingreso, endeudamiento, moras, antigüedad, empleo formal.
\item \textbf{+ Red estática}: de quién y cuánto recibe, concentración del ingreso,
PageRank, salud del empleador y de los pagadores.
\item \textbf{+ Red dinámica}: coherencia con los pagadores, desfase, ritmo,
irregularidad de los ciclos, pagadores que dejaron de pagar.
\end{enumerate}
\end{representa}

\begin{aviso}
Requiere los resultados del experimento. Si no existen, ejecute
\texttt{uv run pulso credit configs/credit.yaml} (menos de un minuto).
\end{aviso}

\subsection*{Control}
\begin{controles}
Modelo & XGBoost o Regresión logística & Modelo con que se comparan los niveles en la
parte superior. La regresión logística es el modelo clásico de los bancos; XGBoost
combina árboles de decisión. \\
\end{controles}

\subsection*{Comparación de niveles}
\begin{description}
\item[Dinámica vs.\ red estática ($\Delta$AUC).] Cuánto mejora el AUC al agregar las
variables dinámicas a las estáticas, en el mundo con red. Debajo, el intervalo de
confianza del 95~\%: si no incluye el cero, la mejora es significativa.
\item[Híbrido vs.\ clásico ($\Delta$AUC).] Mejora total del modelo con todas las variables
sobre el score clásico.
\item[Control nulo.] La misma comparación en el mundo nulo. Debe ser cercana a cero. Si
no lo fuera, habría una fuga de información.
\item[AUC fuera de muestra.] Barras por mundo y nivel. El AUC es la probabilidad de que el
modelo asigne más riesgo a un cliente que cae que a uno que no cae: 0.5 es azar, 1 es
perfecto. Las líneas negras son la variación entre particiones.
\item[PR-AUC.] Métrica similar, más exigente cuando los impagos son pocos. Un modelo al
azar obtiene la tasa de impago (cerca de 0.1).
\item[Mejoras pareadas por partición.] Cada fila es una comparación con su intervalo de
confianza. En verde, las que están por encima de cero.
\item[Tabla completa de métricas.] Desplegable con AUC, PR-AUC, KS (separación entre
buenos y malos pagadores), Brier y ECE (calibración: menor es mejor).
\end{description}

\captura{hibrido_2.png}

\subsection*{Riesgo por cliente en la cartera activa}
Usa la cartera de la página \emph{Cartera multicapa} (personas, firmas, informalidad y
semilla), siempre en el mundo con red y con los valores de $K$ y $\gamma$ del archivo de
configuración, no los de los deslizadores. Las probabilidades salen de modelos XGBoost
finales entrenados con las 20 carteras del experimento.

\begin{description}
\item[PD a 12 meses: score clásico vs.\ híbrido.] Cada punto es un cliente activo. Eje
horizontal: probabilidad de impago según el score clásico. Eje vertical: según el modelo
híbrido. \textcolor{rojo}{Rojo}: el cliente cayó; \textcolor{verde}{verde}: no cayó.
Cuadrados: empresas. Los puntos por encima de la diagonal son clientes que el híbrido ve
más riesgosos que el clásico.
\item[Riesgo oculto.] Los 10 clientes cuyo riesgo más sube al mirar la red, con ambas PD y
si cayeron.
\item[Tasa de impago en el 10~\% más riesgoso.] De los clientes que cada modelo marca
como el 10~\% más riesgoso, cuántos cayeron. Con una sola cartera pequeña esta cifra
varía mucho; la comparación válida es la de la parte superior.
\end{description}

\captura{hibrido_3.png}

\subsection*{¿Por qué? Contribuciones SHAP}
\begin{description}
\item[Importancia media $|$SHAP$|$.] Las 15 variables que más pesan en el modelo híbrido.
Gris: clásicas. Azul: red estática. Rojo: red dinámica.
\item[Explicar un cliente.] Elija un cliente (primero aparecen los de riesgo oculto). Las
barras muestran qué sube (rojo) y qué baja (verde) su riesgo, en log-odds: la escala en
que el modelo suma contribuciones.
\item[\boton{Descargar PD por cliente (CSV)}.] Descarga ambas PD de todos los clientes.
\end{description}

\begin{pruebe}
Cambie el modelo entre XGBoost y regresión logística: la conclusión se mantiene. Después
abra el cliente con mayor riesgo oculto y vea qué variables de red explican el aumento.
\end{pruebe}
